% ======================================================================
\section{Conclusions}\label{sec:conclusions}
% ======================================================================

We have computed the single inclusive gluon spectrum at mid rapidity in the collision of a dilute projectile with a dense target to next-to-leading order in the Color Glass Condensate, in pure Yang--Mills theory. The calculation uses the evolution operator $\Omega$ of the projectile, written as a normal-ordered series in the soft gluon operators with coefficients that are functionals of the non-commuting valence charge density, and reduces the spectrum to products of these coefficients and of Wilson lines. The cross section at order $g^4$ is organised into three groups of terms and fifteen rows. The colour charges of the projectile have been treated as non-commuting operators throughout; their ordering, the contact terms left by their commutators and the symmetrised three-point function of the projectile are part of the result, which is summarised in Section~\ref{sec:result_summary}.

The rapidity-enhanced terms have the structure required by high-energy factorisation. The logarithm of the longitudinal cut-offs splits at the momentum of the measured gluon. The coefficient of $\log(\vee/k^+)$ is the BFKL evolution of the projectile charge correlator inserted into the LO spectrum, and the coefficient of $\log(k^+/\Lambda)$ is the JIMWLK Hamiltonian acting on the target operator of the LO spectrum, for an arbitrary ordered projectile correlator. Terms with four charges cancel for commuting charges, and the terms with three charges that carry a logarithm cancel in the symmetrised sense; with non-commuting charges, their reordering produces part of the two-charge structures of both evolutions. The operator ordering of the charges, in particular in the three-charge term of the one-gluon component of the wave function at order $g^3$, is essential for these cancellations.

All poles $1/\epsilon$ of the calculation are accounted for, in conventional dimensional regularisation with the $\MS$ coupling and the $\MS$ fragmentation function. The ultraviolet poles of the vertex correction of the measured gluon and of the term in which the two gluons of the two-gluon component recombine cancel each other; the finite terms that accompany this cancellation depend on the grouping of the two rows. The coupling is renormalised by the ultraviolet half of the normalisation of the dressed measured gluon, and the infrared half of that normalisation completes the real final-state splitting to the full $\MS$ splitting function, whose collinear pole is absorbed into the fragmentation function. A left-over pole multiplied by $\log(1-k^+/\vee)$ is a power correction outside the eikonal approximation, which we define as the leading term of the expansion in $k^+/\vee$ and $\Lambda/k^+$ at fixed $\epsilon$. The collinear term is the final-state DGLAP fragmentation of the measured gluon, with its end points subtracted into the JIMWLK and the BFKL term. The running-coupling term is the $\frac{\beta_0}{2}\delta(1-\xi)$ part of the collinear term, added and subtracted; with it, the coupling of the LO spectrum runs with the transverse separation conjugate to the measured momentum or, in the local form, with the measured momentum itself.

The finite terms are given Group by Group and Row by Row, for an arbitrary projectile and, after the average over a globally colour-invariant projectile, in colour-singlet form with traces of adjoint Wilson lines. They contain no renormalisation scale and no rapidity cut-off. The soft remainders of the plus prescriptions in the rows in which the measured gluon recoils against a soft gluon cancel in pairs of a virtual and a real term, without any condition on the target at large distance, and the large-distance tail of the collinear row cancels against two other rows. In the averaged form the terms cubic in the charges contribute in two ways: their symmetrised part is proportional to the symmetrised three-point function of the projectile, for which no model is assumed, and is odd in $\v k$; the commutators of the ordered charges leave contact terms of the same order as the quadratic terms, which are even in $\v k$ and present for any projectile. The vanishing energy denominators, which occur in one row, are defined as principal values.

In the limit of a dilute target the real emission reproduces the two-gluon production cross section of quasi-multi-Regge kinematics~\cite{Fadin:1998py,Bartels:2006hg}, including its normalisation and its polarisation structure. For a point-like projectile the virtual correction has the gluon Regge trajectories as coefficients of the two rapidity logarithms; its comparison with the one-loop central-emission vertex has not been completed.

Rapidity evolution of the projectile and of the target, collinear evolution of the final state, the running of the coupling and the finite terms thus emerge from one calculation, which completes the NLO description of the single inclusive gluon spectrum within the eikonal, dilute-projectile approximation, up to the open points listed in Section~\ref{sec:openpoints}. The most important of these are the comparison of the virtual correction with the one-loop vertex in the dilute limit, the disagreement in the parity of the cubic terms with~\cite{Li:2021ntm}, the dependence on the prescription for the vanishing energy denominators, and the numerical evaluation of the finite terms, which has not been carried out.
